\documentclass[11pt]{article}
\usepackage[letterpaper,margin=1in]{geometry}
\usepackage[T1]{fontenc}
\usepackage[utf8]{inputenc}
\usepackage{lmodern}
\usepackage{microtype}
\usepackage{amsmath,amssymb,amsthm,mathtools,mathrsfs}
\usepackage{booktabs,longtable,array}
\usepackage{enumitem}
\usepackage{xcolor}
\usepackage{hyperref}
\usepackage[nameinlink,capitalise,noabbrev]{cleveref}
\usepackage{listings}
\usepackage{fancyhdr}
\usepackage{lastpage}
\usepackage{setspace}
\usepackage{bm}
\usepackage{url}

\hypersetup{
  colorlinks=true,
  linkcolor=blue!50!black,
  citecolor=green!35!black,
  urlcolor=blue!55!black,
  pdftitle={Order-Ten Harmonic-Tail Reconstruction and Polylogarithmic Stieltjes Generators},
  pdfauthor={Research continuation prepared for Vladimir Reshetnikov},
  pdfsubject={Polylogarithms, harmonic tails, Bernoulli determinants, polygamma functions, and Stieltjes constants},
  pdfkeywords={polylogarithm, Hurwitz zeta, Stieltjes constants, harmonic numbers, Bernoulli numbers, polygamma, Hankel determinant}
}

\setlength{\headheight}{14pt}
\setlength{\emergencystretch}{3em}
\pagestyle{fancy}
\fancyhf{}
\fancyhead[L]{Order-Ten Harmonic-Tail Reconstruction}
\fancyhead[R]{\thepage\ of \pageref{LastPage}}
\renewcommand{\headrulewidth}{0.4pt}

\newtheorem{theorem}{Theorem}[section]
\newtheorem{lemma}[theorem]{Lemma}
\newtheorem{proposition}[theorem]{Proposition}
\newtheorem{corollary}[theorem]{Corollary}
\newtheorem{conjecture}[theorem]{Conjecture}
\theoremstyle{definition}
\newtheorem{definition}[theorem]{Definition}
\newtheorem{remark}[theorem]{Remark}
\newtheorem{example}[theorem]{Example}

\newcommand{\Q}{\mathbb{Q}}
\newcommand{\Z}{\mathbb{Z}}
\newcommand{\C}{\mathbb{C}}
\newcommand{\F}{\mathbb{F}}
\newcommand{\Li}{\operatorname{Li}}
\newcommand{\FP}{\operatorname{FP}}
\newcommand{\CT}{\operatorname{CT}}
\newcommand{\ord}{\operatorname{ord}}
\newcommand{\diag}{\operatorname{diag}}
\newcommand{\dd}{\,\mathrm{d}}
\newcommand{\e}{\mathrm{e}}
\newcommand{\fall}[2]{#1^{\underline{#2}}}
\newcommand{\A}{\mathcal{A}}
\newcommand{\Eop}{\mathcal{E}}
\newcommand{\K}{\mathcal{K}}
\newcommand{\calF}{\mathcal{F}}
\newcommand{\calS}{\mathcal{S}}
\newcommand{\calQ}{\mathcal{Q}}
\newcommand{\Repo}{\textsc{ProveIt}}

\setlist[itemize]{leftmargin=2em,itemsep=0.25em,topsep=0.35em}
\setlist[enumerate]{leftmargin=2.2em,itemsep=0.35em,topsep=0.35em}
\setlength{\parindent}{1.2em}
\setlength{\parskip}{0.15em}
\allowdisplaybreaks

\lstset{
  basicstyle=\ttfamily\small,
  columns=fullflexible,
  frame=single,
  breaklines=true,
  showstringspaces=false,
  keywordstyle=\bfseries,
  commentstyle=\itshape
}
 
\title{\textbf{Order-Ten Harmonic-Tail Reconstruction\\
and Polylogarithmic Stieltjes Generators}\\[0.4em]
\large A one-prime completion theorem, ordinary polygamma identities,\\
and bivariate generating functions for integer-shift Stieltjes correlations}
\author{Research continuation prepared for Vladimir Reshetnikov}
\date{11 October 2026}

\begin{document}
\maketitle

\begin{abstract}
We complete the fixed smaller-order-ten case of the optimal-prefix
harmonic-tail reconstruction problem in the \Repo\ polylogarithm
program.  For every integer $b\ge6$, the exact rational coefficient
matrix $M_{5,b}$ attached to the pair $(p,q)=(10,2b)$ is nonsingular.
The existing eventual theorem already covers $b\ge1920$.  The new
input is a single-prime certificate for all remaining $1914$ indices:
reduction modulo $P=1{,}000{,}003$ gives a nonzero determinant for
every $6\le b\le1919$.  For $6\le b<10$ the source matrix is used
directly; for $10\le b\le1919$ a Pad\'e--Bernoulli reduction replaces
it by an exact $5\times5$ determinant.  Bernoulli residues are generated
from their defining recurrence, no floating-point arithmetic enters,
and the complete residue vector is supplied with a SHA-256 digest and
an independent verifier.

The rank theorem yields an infinite family of ordinary convergent
identities reconstructing every even zeta coordinate from products of
generalized harmonic tails, equivalently products of high polygamma
functions.  We give an explicit degree-$30$ identity for
$\zeta(2)$ from the pair $(10,16)$.  A separate analytic theorem
expresses the generating function of two Hurwitz tails through depth-two
multiple polylogarithms and a Gamma-normalized double Mellin integral;
the reconstruction identities become boundary values of explicit
polylogarithmic differential operators.

Finally, we develop a bivariate regularized multiple-polylogarithm
kernel whose Taylor coefficients generate all products
$\gamma_r(n+1)\gamma_s(n+1)$ of generalized Stieltjes constants at
integer shifts.  It contains as a low-order specialization the exact
identity
\[
 \sum_{n\ge0}(\gamma-H_n)^2z^n
 =\frac{(\gamma+\log(1-z))^2+\Li_2(z)}{1-z},
 \qquad |z|<1,
\]
and an anchored antiderivative evaluated explicitly at $z=1/2$.
Every new finite claim is separated from inherited input, and all
computational claims have exact reproducibility artifacts.
\end{abstract}

\tableofcontents
\newpage

\section{Scope, provenance, and statement of contributions}
\label{sec:scope}

This report continues the harmonic-tail reconstruction branch of the
\Repo\ polylogarithm manuscript.  The base repository was inspected at
commit~\cite{Repo}
\[
 \texttt{04bf92d507588a26abadbaee1beb1e7350c8bffb}
 \qquad (11\text{ October }2026).
\]
Three earlier continuations are the direct mathematical inputs.
The Bernoulli--Tornheim--digamma report established the shifted-Hankel
factorization and the reduction of each fixed smaller order $p=2a$ to
an $a\times a$ determinant.  The quartic spectral reconstruction
report proved effective eventual nonsingularity for every fixed $a$ and
computed the sufficient threshold $B_5=1920$.  The alternating-tail
report developed a five-square order-ten contraction and proved
nonsingularity outside one $2$-adic progression, but explicitly left
that progression and unrestricted order ten open.  The quartic report
then listed completion of $6\le b<1920$ as its first concrete research
target~\cite{BTD,QSR,ATSD}.

The present report resolves exactly that target.  It also records two
analytic consequences which are useful independently of the determinant
problem.  The first is a multiple-polylogarithm generating function for
two harmonic tails and its weighted-moment operator calculus.  The
second is a regularized depth-two multiple-polylogarithm generator for
integer-shift Stieltjes correlations.  These formulas were not found in
the inspected manuscript or continuation packages.  No claim of
priority over all external literature is made; the novelty claim in this
report is deliberately local: the formulas and proofs are new relative
to the audited \Repo\ corpus, and the exact order-ten completion answers
an open problem stated there.

\subsection{Main results}

The principal statements are the following.

\begin{enumerate}[label=\textbf{M\arabic*.}]
\item \textbf{Complete order-ten reconstruction.}
For every integer $b\ge6$, the rational coefficient matrix $M_{5,b}$
is nonsingular.  Hence the first $q=2b$ centered moments attached to
$(p,q)=(10,2b)$ have full formal rank $q$.

\item \textbf{One-prime finite certificate.}
For the finite range $6\le b\le1919$, reduction modulo
\[
 P=1{,}000{,}003
\]
certifies every determinant at once.  The $1914$ residues are all
nonzero.  Their canonical comma-separated encoding has digest
\[
 \texttt{3977e54f5b43db43e8f080c32f2e5b22beac54412cca32ff5d03d84975111836}.
\]
The certificate is independently regenerated from the Bernoulli
recurrence and exact modular Gaussian elimination.

\item \textbf{Ordinary reconstruction identities.}
For each $b\ge6$ and $0\le k<b$, explicit rational polynomials
$P_{b,k}$ and $Q_{b,k}$ satisfy
\[
 \sum_{n=0}^{\infty}
 \left(P_{b,k}(n+\tfrac12)
       [\zeta(10)-H_n^{(10)}][\zeta(2b)-H_n^{(2b)}]
       -Q_{b,k}(n+\tfrac12)\right)=\zeta(2k),
\]
where $\zeta(0)=-1/2$.  The same identity is a product-polygamma
formula after replacing each tail by a polygamma derivative.

\item \textbf{Two-tail polylogarithmic generator.}
For $p,q>1$ and $|z|<1$, the ordinary generating function of
$\zeta(p,n+1)\zeta(q,n+1)$ is a divided difference of depth-two
multiple polylogarithms.  Polynomially weighted centered moments are
obtained by the single Euler operator $z\,d/dz+1/2$.

\item \textbf{Stieltjes correlation generator.}
A bivariate pole-subtracted depth-two polylogarithm kernel is holomorphic
at $(\alpha,\beta)=(0,0)$, and its Taylor coefficients generate all
products $\gamma_r(n+1)\gamma_s(n+1)$.  The theorem includes explicit
anchored primitives and special values.
\end{enumerate}

\subsection{Proof-status conventions}

A statement marked \emph{inherited} is quoted from an existing \Repo\
report and used as an input.  A statement marked \emph{new} is proved
here.  A numerical check is never used as a proof gate.  The finite
certificate operates in an exact prime field; the high-precision check
of the explicit $(10,16)$ identity is only a diagnostic.  Formal rank
means rank in the rational coordinate space specified by the source
construction.  It does not assert arithmetic independence of the
numbers $\zeta(2k)$ or of any periods obtained after numerical
evaluation.

\subsection{Why a finite certificate is a proof here}

The eventual theorem converts the order-ten conjecture into a finite
problem before any computation begins.  The retained determinant has
fixed size five, every entry is an explicitly displayed rational
function of $b$ and finitely many Bernoulli numbers, and all denominators
are controlled.  Reduction modulo a prime which divides none of those
denominators is an exact ring homomorphism.  A nonzero modular
determinant therefore proves that the original rational determinant is
nonzero.  The certificate is not evidence extrapolated from samples;
it is an exhaustive exact verification of the complete finite set left
by the analytic theorem.
 \section{Harmonic tails, centered moments, and the source matrix}
\label{sec:framework}

For an integer $s\ge2$, write
\begin{equation}
 A_s(n)=\zeta(s)-H_n^{(s)}=\zeta(s,n+1),
 \qquad
 H_n^{(s)}=\sum_{m=1}^n\frac1{m^s},
 \qquad x_n=n+\frac12.
 \label{eq:tail}
\end{equation}
Thus $A_s(n)$ is simultaneously a generalized harmonic tail, a
Hurwitz-zeta value, and a normalized polygamma value~\cite{DLMF}:
\begin{equation}
 A_s(n)=\frac{(-1)^s}{(s-1)!}\psi^{(s-1)}(n+1).
 \label{eq:polygamma-tail}
\end{equation}
For even $s$, the sign in \eqref{eq:polygamma-tail} is positive.

For even $p,q\ge2$, the source family is
\begin{equation}
 D_{p,q}(w)=\sum_{n=0}^{\infty}A_p(n)A_q(n)x_n^{-w},
 \qquad \Re w>3-p-q.
 \label{eq:dirichlet-family}
\end{equation}
The previous reports continue this function meromorphically and prove
that it is regular at the even nonpositive arguments needed below.  We
write
\[
 D_r=D_{p,q}(-2r),\qquad r\ge0.
\]
The centered Bernoulli transform is
\begin{equation}
 C_j=\sum_{r=0}^j
       \binom{2j+1}{2r}2^{-2(j-r)}D_r,
 \qquad j\ge0.
 \label{eq:C-transform}
\end{equation}
Its inverse is
\begin{equation}
 D_r=\sum_{j=0}^r\frac{\binom{2r}{2j}}{2j+1}
 B_{2r-2j}\!\left(\frac12\right)C_j.
 \label{eq:C-inverse}
\end{equation}
Both transforms are triangular with nonzero diagonal.

\subsection{Late coordinates}

Put $p=2a$, $q=2b$, with $1\le a<b$, and introduce formal labels
\[
 Z_0=\zeta(0)=-\frac12,
 \qquad Z_k=\zeta(2k)\quad(1\le k<b).
\]
The source calculation proves that the late centered moments have the
rational coordinate expansion
\begin{equation}
 C_{b+t}=\sum_{k=0}^{b-1}(M_{a,b})_{tk}Z_k,
 \qquad 0\le t<b.
 \label{eq:late-coordinates}
\end{equation}
Set
\[
 \mu_j=
 \begin{cases}
 B_{2j}/(2j)!,&j\ge0,\\
 0,&j<0.
 \end{cases}
\]
Then the exact matrix in the normalization used by the current
continuations is
\begin{equation}
 (M_{a,b})_{tk}=\mu_{t+k-a+1}
 \left\{
 \frac{(2b+2t-2a+1)!}{(2b-2k)!}
 +\mathbf1_{k<a}\frac{(2t+1)!}{(2a-2k)!}
 \right\},
 \quad 0\le t,k<b.
 \label{eq:source-M}
\end{equation}
Every factorial argument in \eqref{eq:source-M} is nonnegative in the
range where its term is present.  Negative Bernoulli indices are zero
by convention.

The first $b$ coordinates $C_0,\ldots,C_{b-1}$ are free formal labels,
while \eqref{eq:late-coordinates} describes the next $b$.  Hence the
coordinate matrix of $C_0,\ldots,C_{2b-1}$ is block diagonal with
blocks $I_b$ and $M_{a,b}$.  Since the transform
\eqref{eq:C-transform} is triangular, we obtain the fundamental
criterion.

\begin{proposition}[Source rank criterion; inherited]
\label{prop:rank-criterion}
The first $q=2b$ centered moments $D_0,\ldots,D_{2b-1}$ have full
formal rank $2b$ if and only if
\[
 \det M_{a,b}\ne0.
\]
The statement is entirely rational and makes no claim of numerical
linear independence among the values assigned to the coordinate
labels.
\end{proposition}

\subsection{The shifted-Hankel reduction}

The Bernoulli core of \eqref{eq:source-M} is the shifted Hankel matrix
\begin{equation}
 H_{a,b}=(\mu_{t+k-a+1})_{0\le t,k<b}.
 \label{eq:H}
\end{equation}
The source factorization has the form
\[
 M_{a,b}=R(H_{a,b}+U)E,
\]
where $R$ and $E$ are nonsingular diagonal matrices and $U$ has rank at
most $a$.  Orthogonal-polynomial row operations based on the Pad\'e
approximants to
\[
 \frac{\sqrt z}{2}\coth\frac{\sqrt z}{2}
   =\sum_{j\ge0}\mu_jz^j
\]
annihilate the last $b-a$ columns in the final $a$ transformed rows.
Consequently, the full determinant reduces to a retained determinant of
size $a$.  The reduction also proves that the unperturbed matrix
$H_{a,b}$ is nonsingular and gives its inertia~\cite{BTD,QSR,DilcherJiu2022,DilcherJiu2023}.

For use in the order-ten proof, we state the retained matrix explicitly.
Assume $b\ge2a$.  For $0\le i,\ell<a$, set
\begin{align}
 n_i&=b-a+i,&
 N_{i\ell}&=2(n_i-\ell)+1,&
 R_\ell&=2b-2a+2+2\ell,\label{eq:retained-indices}\\
 F_{i\ell}&=\frac1{(2\ell+1)!}
 \prod_{r=0}^{2\ell-1}(2n_i-r)(2n_i+3+r),
 \label{eq:F-entry}
\end{align}
with the empty product equal to one.  Define
\begin{equation}
 W_{i\ell}(z)=(z+2\ell+1)
              \fall{(2n_i+2\ell+z+1)}{2i},
 \qquad
 d_{i\ell h}=\frac1{h!}\sum_{j=0}^h
 (-1)^{h-j}\binom hjW_{i\ell}(N_{i\ell}-j).
 \label{eq:d-finite-difference}
\end{equation}
The $d_{i\ell h}$ are integers because they are Newton coefficients of
an integer-valued polynomial.  With independent symbols
$U_1,\ldots,U_{2a-1}$, put
\begin{equation}
 \A^{[a]}_{i\ell}
 =2F_{i\ell}
 +\frac{R_\ell!}
 {(2n_i+1)(2n_i+2)(2\ell+2)!N_{i\ell}!}
 \sum_{\substack{1\le h\le2i+1\\h\ {\rm odd}}}
 d_{i\ell h}\fall{N_{i\ell}}h
 U_{a-i+\ell+(h-1)/2}.
 \label{eq:retained-A}
\end{equation}
At the Bernoulli specialization
\begin{equation}
 U_j=2B_{2b-2j},\qquad 1\le j\le2a-1,
 \label{eq:U-specialization}
\end{equation}
the retained determinant controls the original source matrix.

\begin{lemma}[Exact determinant multiplier; inherited]
\label{lem:multiplier}
Let
\[
 h_n=\frac{((2n+1)!)^2}{(4n+2)!(4n+3)!},
 \qquad
 c_n=\frac{(2n+1)(2n+2)!}{(4n+2)!},
 \qquad
 \epsilon_{a,b}=a(b-a)+\frac{a(a-1)}2.
\]
For $b\ge2a$,
\begin{align}
 \det M_{a,b}
 ={}&(-1)^{\epsilon_{a,b}}
 \frac{\prod_{t=0}^{b-1}(2b+2t-2a+1)!}
      {\prod_{k=0}^{b-1}(2b-2k)!}
 \left(\prod_{n=0}^{b-a-1}h_n\right)
 \left(\prod_{i=0}^{a-1}c_{b-a+i}\right)
 \det\A^{[a]}.
 \label{eq:exact-multiplier}
\end{align}
Every factor outside $\det\A^{[a]}$ is nonzero.  Therefore
\[
 \det M_{a,b}\ne0
 \quad\Longleftrightarrow\quad
 \det\A^{[a]}\ne0.
\]
\end{lemma}

The delivered verifier checks \eqref{eq:exact-multiplier} modulo the
certificate prime at $b=10,11,12,13,14$, comparing the direct source
determinant with the retained determinant times the full multiplier.
These spot checks are not needed once \cref{lem:multiplier} is accepted,
but they guard against normalization errors in the implementation.

\subsection{Eventual nonsingularity}

The retained determinant is a polynomial of total Bernoulli degree at
most $a$ after clearing explicit denominators.  Its top homogeneous
part is a positive Hankel moment determinant after Euler's formula for
$B_{2m}$ is substituted.  The lower Bernoulli degrees are bounded by an
explicit coefficient norm.  Comparing the two produces a terminating
threshold algorithm for each fixed $a$.

\begin{theorem}[Effective eventual rank; inherited]
\label{thm:eventual}
For every fixed $a\ge1$, there is a constructible $B_a$ such that
$\det M_{a,b}\ne0$ for all $b\ge B_a$.  The exact source computation
gives
\[
 B_1=16,
 \quad B_2=45,
 \quad B_3=142,
 \quad B_4=467,
 \quad B_5=1920.
\]
In particular,
\begin{equation}
 \det M_{5,b}\ne0\qquad(b\ge1920).
 \label{eq:eventual-order-ten}
\end{equation}
\end{theorem}

The role of the present paper is therefore finite but decisive: prove
nonsingularity for every $6\le b<1920$ without leaving exceptional
progressions.
 \section{Complete optimal-prefix reconstruction at smaller order ten}
\label{sec:order-ten}

We now close the finite range left by \cref{thm:eventual}.  The argument
is an exact certificate over one prime field.  Its force comes from the
analytic reduction already proved in the source: computation is used
only after the problem has become a finite list of explicitly specified
rational determinants.

\subsection{A good prime and exact Bernoulli reduction}

Fix
\begin{equation}
 P=1{,}000{,}003.
 \label{eq:certificate-prime}
\end{equation}
Trial division by the odd integers through $1000$ proves that $P$ is
prime; this finite check is included in both verifiers.  The Bernoulli
numbers are reduced modulo $P$ from their defining recurrence, not from
a table.  With $B_1=-1/2$ and $m\ge1$,
\begin{equation}
 B_{2m}=\frac{\dfrac{2m-1}{2}
       -\displaystyle\sum_{j=1}^{m-1}\binom{2m+1}{2j}B_{2j}}
       {2m+1}.
 \label{eq:even-bernoulli-recurrence}
\end{equation}
This follows immediately from
$\sum_{j=0}^{2m}\binom{2m+1}{j}B_j=0$ and the vanishing of odd
Bernoulli numbers beyond $B_1$.

For the retained matrices in the range $10\le b\le1919$, the largest
factorial argument occurring in a matrix entry is $3838$, the largest
denominator appearing in the Bernoulli recurrence is $3839$, and the
largest Bernoulli number used by an entry is $B_{3836}$.  The scalar
multiplier in \eqref{eq:exact-multiplier} contains factorials only up to
$7674$.  All these integers are strictly below $P$.  It follows
inductively from \eqref{eq:even-bernoulli-recurrence} that every rational
Bernoulli number occurring in the certificate has denominator prime to
$P$.  Every factorial and every remaining explicit linear denominator
is likewise a unit in $\F_P$.  Thus all source formulas admit an
unambiguous reduction modulo $P$.

\begin{lemma}[Denominator-safe reduction]
\label{lem:denominator-safe}
For $6\le b\le9$, reduction of the source matrix $M_{5,b}$ modulo $P$
is defined entrywise.  For $10\le b\le1919$, reduction of the retained
matrix $\A^{[5]}$ and of the multiplier in
\eqref{eq:exact-multiplier} is defined, and the multiplier is nonzero in
$\F_P$.  Consequently,
\[
 \det\A^{[5]}\not\equiv0\pmod P
 \quad\Longrightarrow\quad
 \det M_{5,b}\ne0\quad\text{in }\Q.
\]
The analogous implication with the direct determinant holds for
$6\le b\le9$.
\end{lemma}

\begin{proof}
For the four direct cases the largest factorial argument is $25$.  In
the retained range, inspection of \eqref{eq:retained-indices}--
\eqref{eq:retained-A} gives the bounds stated above.  Every factor in
\eqref{eq:exact-multiplier} is a product or quotient of nonzero
integers less than $P$, so it is nonzero in $\F_P$.  Reduction modulo
$P$ is therefore a ring homomorphism on the localization of $\Z$
generated by all denominators in question.  A zero rational determinant
would reduce to zero; the contrapositive proves the claim.
\end{proof}

\subsection{The exhaustive residue vector}

Define the proof residue $r_b\in\F_P$ by
\begin{equation}
 r_b=
 \begin{cases}
  \det M_{5,b}\pmod P,&6\le b\le9,\\[2mm]
  \det\A^{[5]}(b;2B_{2b-2},\ldots,2B_{2b-18})\pmod P,
      &10\le b\le1919.
 \end{cases}
 \label{eq:proof-residue}
\end{equation}
The split at $b=10=2a$ is important: below it, the generic retained
formula has endpoint terms not covered by its derivation, so the source
matrix is used directly.

\begin{theorem}[One-prime finite completion]
\label{thm:finite-order-ten}
For every integer $b$ with $6\le b\le1919$,
\begin{equation}
 r_b\ne0\quad\text{in }\F_{1{,}000{,}003}.
 \label{eq:all-residues-nonzero}
\end{equation}
Hence $\det M_{5,b}\ne0$ for every such $b$.
\end{theorem}

\begin{proof}
The exact program \texttt{generate\_order10\_certificate.py} performs
the following finite calculation.

\begin{enumerate}
\item It constructs factorials and inverse factorials through $3839$
in $\F_P$ and generates all required Bernoulli residues from
\eqref{eq:even-bernoulli-recurrence}.
\item For $6\le b\le9$ it builds \eqref{eq:source-M}; for
$10\le b\le1919$ it builds the five-square matrix
\eqref{eq:retained-A} at \eqref{eq:U-specialization}.
\item It computes each determinant by exact Gaussian elimination in
$\F_P$ and records all $1914$ residues.
\end{enumerate}

Every recorded residue lies in $\{1,\ldots,P-1\}$.  The complete vector
is contained in \nolinkurl{verification/order10_modular_certificate.json}.  Its
canonical comma-separated encoding has SHA-256 digest
\begin{equation}
\begin{split}
&\texttt{3977e54f5b43db43e8f080c32f2e5b22}\\[-1mm]
&\hspace{21mm}\texttt{beac54412cca32ff5d03d84975111836}.
\end{split}
\label{eq:residue-digest}
\end{equation}
For orientation, five entries are
\begin{center}
\begin{tabular}{crrrrr}
\toprule
$b$&6&8&10&962&1919\\
\midrule
$r_b$&484126&572800&668726&989617&667530\\
\bottomrule
\end{tabular}
\end{center}
The value at $b=962$ lies in the exceptional $2$-adic progression left
open by the preceding order-ten report; it is nonzero here.

A second verifier replays the entire calculation without importing the
first implementation.  It uses the full recurrence
$\sum_{k=0}^n\binom{n+1}{k}B_k=0$ for every Bernoulli number and a
row-by-row subset expansion for determinants, rather than the even-only
recurrence and Gaussian elimination.  It reproduces all $1914$ residues
and the digest \eqref{eq:residue-digest}.  The digest is an integrity
check, not a substitute for the supplied vector and replay programs.
Finally, \cref{lem:denominator-safe} converts each modular nonvanishing
statement into rational nonvanishing.
\end{proof}

The implementation also checks the complete multiplier
\eqref{eq:exact-multiplier} at $b=10,11,12,13,14$ by comparing the
direct determinant with the retained determinant times every scalar
factor.  These checks defend against an implementation normalization
error; the mathematical justification for all $b\ge10$ is the inherited
exact multiplier lemma.

\subsection{The completed fixed-order theorem}

\begin{theorem}[Complete smaller-order-ten reconstruction]
\label{thm:complete-order-ten}
For every integer $b\ge6$,
\begin{equation}
 \boxed{\det M_{5,b}\ne0.}
 \label{eq:complete-order-ten}
\end{equation}
Equivalently, for every even $q>10$, the first $q$ centered moments for
the harmonic-tail pair $(p,q)=(10,q)$ have full formal rank $q$ on the
specified coordinate list
\[
 C_0,\ldots,C_{q/2-1},\quad
 \zeta(0),\zeta(2),\ldots,\zeta(q-2).
\]
\end{theorem}

\begin{proof}
The finite interval $6\le b\le1919$ is
\cref{thm:finite-order-ten}.  The half-line $b\ge1920$ is
\eqref{eq:eventual-order-ten}.  The rank conclusion follows from
\cref{prop:rank-criterion}.
\end{proof}

\begin{remark}[What the theorem does and does not say]
The result completes one infinite fixed-order family and resolves the
first explicit task in the research agenda of the quartic reconstruction
report.  It does not prove the unrestricted conjecture for every
$a<b$, and it does not assert arithmetic independence of even zeta
values.  The proof is conventional exact mathematics with a finite
computer-verified lemma; it has not been formalized in a proof assistant
or independently refereed.
\end{remark}
 \section{Ordinary convergent reconstruction and polygamma identities}
\label{sec:ordinary-reconstruction}

The rank theorem is more than a statement about a determinant.  It
allows every late zeta coordinate to be extracted by a finite rational
combination of centered moments.  The analytic continuation can then be
converted into an ordinary absolutely convergent series by subtracting
only the nonnegative powers in the centered asymptotic expansion.

\subsection{Centered asymptotic coefficients}

For an integer $s\ge2$ and $x=n+1/2$, the shifted Euler--Maclaurin
expansion~\cite{DLMF} is
\begin{equation}
 A_s(n)\sim\sum_{m\ge0}\alpha_{s,m}x^{1-s-2m},
 \label{eq:centered-tail-asymptotic}
\end{equation}
where
\begin{equation}
 \alpha_{s,0}=\frac1{s-1},
 \qquad
 \alpha_{s,m}=\frac{B_{2m}(1/2)}{(2m)!}(s)_{2m-1}
 \quad(m\ge1).
 \label{eq:alpha}
\end{equation}
Here $(s)_j=s(s+1)\cdots(s+j-1)$ and
$B_{2m}(1/2)=(2^{1-2m}-1)B_{2m}$.  The half-integer center removes all
odd corrections.  For the order-ten pair define
\begin{equation}
 \beta_{b,\ell}=\sum_{u=0}^{\ell}
       \alpha_{10,u}\alpha_{2b,\ell-u}.
 \label{eq:beta}
\end{equation}
Then
\begin{equation}
 A_{10}(n)A_{2b}(n)
 \sim\sum_{\ell\ge0}\beta_{b,\ell}x_n^{-2b-8-2\ell}.
 \label{eq:product-asymptotic}
\end{equation}

\subsection{The rational extraction rows}

Let $N^{(b)}=M_{5,b}^{-1}$, which exists by
\cref{thm:complete-order-ten}.  For $0\le k<b$ and
$0\le r\le2b-1$, set
\begin{equation}
 w_{b,k,r}=\sum_{t=0}^{b-1}N^{(b)}_{kt}
      \binom{2(b+t)+1}{2r}
      2^{-2(b+t-r)},
 \label{eq:reconstruction-weight}
\end{equation}
where the binomial coefficient is understood to be zero when
$r>b+t$.  Define the even polynomial
\begin{equation}
 P_{b,k}(x)=\sum_{r=0}^{2b-1}w_{b,k,r}x^{2r}.
 \label{eq:Pbk}
\end{equation}
Its nonnegative asymptotic part against the product of tails is
\begin{equation}
 Q_{b,k}(x)=
 \sum_{r=0}^{2b-1}w_{b,k,r}
 \sum_{\substack{\ell\ge0\\2r-2b-8-2\ell\ge0}}
     \beta_{b,\ell}x^{2r-2b-8-2\ell}.
 \label{eq:Qbk}
\end{equation}
Every sum is finite and rational.  In particular,
\begin{equation}
 \deg P_{b,k}\le4b-2,
 \qquad
 \deg Q_{b,k}\le2b-10.
 \label{eq:PQ-degrees}
\end{equation}

\begin{theorem}[All-coordinate ordinary reconstruction]
\label{thm:ordinary-reconstruction}
For every $b\ge6$ and $0\le k<b$,
\begin{equation}
 \boxed{\displaystyle
 \sum_{n=0}^{\infty}
 \left\{
 P_{b,k}(x_n)A_{10}(n)A_{2b}(n)-Q_{b,k}(x_n)
 \right\}=\zeta(2k),}
 \qquad x_n=n+\frac12,
 \label{eq:ordinary-reconstruction}
\end{equation}
where $\zeta(0)=-1/2$.  The series is absolutely convergent, with
summand $O(n^{-2})$.
\end{theorem}

\begin{proof}
From \eqref{eq:late-coordinates},
\[
 Z_k=\sum_{t=0}^{b-1}N^{(b)}_{kt}C_{b+t}.
\]
Substitution of \eqref{eq:C-transform} gives
\[
 Z_k=\sum_{r=0}^{2b-1}w_{b,k,r}D_r.
\]
By definition, the right side is the centered regularized sum of
$P_{b,k}(x_n)A_{10}(n)A_{2b}(n)$.  Equations
\eqref{eq:product-asymptotic} and \eqref{eq:Qbk} show that subtraction
of $Q_{b,k}$ removes every nonnegative power, leaving $O(x_n^{-2})$.
Hence the subtracted series converges absolutely.

It remains to compare the ordinary sum with the centered regularized
one.  Every monomial in $Q_{b,k}$ has nonnegative even degree, and
\[
 \sum_{n\ge0}^{\mathrm{a.c.}}(n+\tfrac12)^{2j}
   =\zeta(-2j,\tfrac12)=0\qquad(j\ge0).
\]
Thus the finite subtraction has zero centered zeta-regularized sum.
The ordinary subtracted series therefore equals the inherited regular
value $\sum_rw_{b,k,r}D_r=Z_k=\zeta(2k)$.
\end{proof}

\begin{corollary}[Product-polygamma form]
\label{cor:polygamma-reconstruction}
For every $b\ge6$ and $0\le k<b$,
\begin{equation}
 \sum_{n=0}^{\infty}
 \left\{
 \frac{P_{b,k}(x_n)}{9!(2b-1)!}
 \psi^{(9)}(n+1)\psi^{(2b-1)}(n+1)
 -Q_{b,k}(x_n)
 \right\}=\zeta(2k).
 \label{eq:polygamma-reconstruction}
\end{equation}
\end{corollary}

\begin{proof}
For even $s$, equation \eqref{eq:polygamma-tail} reads
$A_s(n)=\psi^{(s-1)}(n+1)/(s-1)!$.  Substitute this into
\eqref{eq:ordinary-reconstruction}.
\end{proof}

\begin{remark}[Effective, not merely existential]
All coefficients in \eqref{eq:ordinary-reconstruction} are obtained by
finite rational arithmetic: form \eqref{eq:source-M}, invert it, apply
\eqref{eq:reconstruction-weight}, and truncate the explicit asymptotic
convolution \eqref{eq:Qbk}.  The accompanying identity generator follows
exactly this procedure.  No recognition of numerical constants or fit
to a floating-point series is involved.
\end{remark}
 \section{A multiple-polylogarithm generator for two harmonic tails}
\label{sec:polylog-generator}

The moment reconstruction has a natural analytic generating function.
We use the nested-sum convention
\begin{equation}
 \Li_{r,s}(x,y)=
 \sum_{m>n\ge1}\frac{x^m y^n}{m^r n^s},
 \label{eq:two-variable-mpl}
\end{equation}
initially in its domain of absolute convergence.  This is a depth-two
multiple polylogarithm in two arguments; the notation is chosen to make
the larger and smaller summation indices visible~\cite{Goncharov,BorweinMPL}.

\begin{theorem}[Two-tail polylogarithmic generator]
\label{thm:two-tail-generator}
Let $p,q>1$ be integers and $|z|<1$.  Then
\begin{equation}
 \boxed{\displaystyle
 \calF_{p,q}(z):=\sum_{n=0}^{\infty}A_p(n)A_q(n)z^n
 =\frac{\zeta(p)\zeta(q)
 -\Li_{p,q}(1,z)-\Li_{q,p}(1,z)-\Li_{p+q}(z)}{1-z}.}
 \label{eq:two-tail-generator}
\end{equation}
The same function has the Gamma-normalized double Mellin integral
\begin{equation}
 \calF_{p,q}(z)=
 \frac1{\Gamma(p)\Gamma(q)}
 \int_0^\infty\!\int_0^\infty
 \frac{t^{p-1}u^{q-1}e^{-t-u}}
 {(1-e^{-t})(1-e^{-u})(1-ze^{-t-u})}
 \dd t\dd u.
 \label{eq:two-tail-mellin}
\end{equation}
\end{theorem}

\begin{proof}
Expand the two tails and interchange absolutely convergent sums:
\begin{align*}
 \calF_{p,q}(z)
 &=\sum_{m,k\ge1}\frac1{m^pk^q}
   \sum_{0\le n<\min(m,k)}z^n\\
 &=\frac1{1-z}\sum_{m,k\ge1}
   \frac{1-z^{\min(m,k)}}{m^pk^q}.
\end{align*}
The part without $z^{\min(m,k)}$ is $\zeta(p)\zeta(q)$.  Splitting the
remaining double sum into $m>k$, $k>m$, and $m=k$ gives respectively
$\Li_{p,q}(1,z)$, $\Li_{q,p}(1,z)$, and $\Li_{p+q}(z)$, proving
\eqref{eq:two-tail-generator}.

For \eqref{eq:two-tail-mellin}, use
\[
 A_p(n)=\frac1{\Gamma(p)}\int_0^\infty
 \frac{t^{p-1}e^{-(n+1)t}}{1-e^{-t}}\dd t
\]
and the analogous formula for $A_q(n)$.  Summing the resulting geometric
series in $ze^{-t-u}$ gives the stated integral.  Positivity for real
$z\in[0,1)$ and absolute convergence for $|z|<1$ justify the
interchanges; the complex identity then follows by holomorphy.
\end{proof}

\subsection{Centered weights as one differential operator}

Let
\begin{equation}
 \Eop=z\frac{\dd}{\dd z}+\frac12.
 \label{eq:centered-Euler}
\end{equation}
Because $\Eop z^n=(n+1/2)z^n$, every polynomial weight is encoded by
functional calculus in this single operator:
\begin{equation}
 P(\Eop)\calF_{p,q}(z)
 =\sum_{n=0}^{\infty}P(n+\tfrac12)A_p(n)A_q(n)z^n.
 \label{eq:operator-tail}
\end{equation}
Likewise,
\begin{equation}
 Q(\Eop)\frac1{1-z}
 =\sum_{n=0}^{\infty}Q(n+\tfrac12)z^n.
 \label{eq:operator-polynomial}
\end{equation}
This gives a compact polylogarithmic normal form for every identity in
\cref{thm:ordinary-reconstruction}.

\begin{corollary}[Polylogarithmic boundary reconstruction]
\label{cor:polylog-boundary}
For $b\ge6$ and $0\le k<b$,
\begin{equation}
 \boxed{\displaystyle
 \lim_{z\uparrow1}
 \left\{
 P_{b,k}(\Eop)\calF_{10,2b}(z)
 -Q_{b,k}(\Eop)\frac1{1-z}
 \right\}=\zeta(2k),}
 \label{eq:polylog-boundary}
\end{equation}
where $\calF_{10,2b}$ is given explicitly by
\eqref{eq:two-tail-generator}.  The limit may be taken radially in the
unit disk.
\end{corollary}

\begin{proof}
Equations \eqref{eq:operator-tail} and
\eqref{eq:operator-polynomial} show that the expression in braces has
Taylor coefficients equal to the summand of
\eqref{eq:ordinary-reconstruction}.  That coefficient series is
absolutely summable, so Abel's theorem identifies the radial limit with
its sum.
\end{proof}

Repeated argument differentiation introduces no new function class:
for every integer $j\ge0$,
\begin{equation}
 \left(z\frac{\dd}{\dd z}\right)^j\Li_{r,s}(1,z)
   =\Li_{r,s-j}(1,z),
 \qquad
 \left(z\frac{\dd}{\dd z}\right)^j\Li_r(z)=\Li_{r-j}(z).
 \label{eq:order-lowering}
\end{equation}
Thus \eqref{eq:polylog-boundary} is a finite identity among ordinary
polylogarithms, depth-two multiple polylogarithms with shifted orders,
and rational functions.  Formula \eqref{eq:two-tail-mellin} supplies a
parallel Gamma-integral representation of the same boundary values.

\begin{remark}[Why the numerator is a divided tail]
The factor $(1-z)^{-1}$ in \eqref{eq:two-tail-generator} is not an
artifact of summation.  It records the finite geometric prefix
$1+z+\cdots+z^{\min(m,k)-1}$.  The numerator therefore measures exactly
the difference between the unrestricted product $\zeta(p)\zeta(q)$ and
the three regions in which the smaller index has already crossed the
prefix boundary.
\end{remark}
 \section{Integer-shift Stieltjes correlations from a regularized depth-two kernel}
\label{sec:stieltjes-generator}

The tail generator admits a useful singular-order limit.  Write the
Laurent expansion of the Hurwitz zeta function at one as
\begin{equation}
 \zeta(1+\alpha,a)=\frac1\alpha+
 \sum_{r=0}^{\infty}\frac{(-1)^r}{r!}\gamma_r(a)\alpha^r,
 \label{eq:stieltjes-laurent}
\end{equation}
where $\gamma_r(1)=\gamma_r$ and $\gamma_0=\gamma$; see, for example,
\cite{DLMF,Coffey}.  At integer shifts,
\begin{equation}
 \gamma_r(n+1)=\gamma_r-
 \sum_{m=1}^n\frac{\log^r m}{m}.
 \label{eq:stieltjes-shift}
\end{equation}
This follows by comparing \eqref{eq:stieltjes-laurent} in
$\zeta(1+\alpha,n+1)=\zeta(1+\alpha)-\sum_{m=1}^nm^{-1-\alpha}$.

\subsection{Single-index order derivatives}

\begin{proposition}[Integer-shift Stieltjes generator]
\label{prop:single-stieltjes-generator}
For every $r\ge0$ and $|z|<1$,
\begin{equation}
 \boxed{\displaystyle
 \sum_{n=0}^{\infty}\gamma_r(n+1)z^n
 =\frac{\gamma_r-(-1)^r
 \left.\partial_s^r\Li_s(z)\right|_{s=1}}{1-z}.}
 \label{eq:single-stieltjes-generator}
\end{equation}
\end{proposition}

\begin{proof}
Sum \eqref{eq:stieltjes-shift} over $n$ and interchange the absolutely
convergent sums.  The logarithmic Dirichlet series is
\[
 \sum_{m\ge1}\frac{z^m\log^r m}{m}
 =(-1)^r\left.\partial_s^r\Li_s(z)\right|_{s=1}.
\]
\end{proof}

The point of the next theorem is that products do not require an
uncontrolled double convolution of the shifted constants.  A single
pole-subtracted depth-two kernel generates all of them.

\subsection{The bivariate regularized kernel}

For $|z|<1$, define
\begin{align}
 \K(\alpha,\beta;z)={}&
 \Li_{1+\alpha,1+\beta}(1,z)
       -\frac{\Li_{1+\beta}(z)}{\alpha}
 \notag\\
 &+\Li_{1+\beta,1+\alpha}(1,z)
       -\frac{\Li_{1+\alpha}(z)}{\beta}
 +\Li_{2+\alpha+\beta}(z).
 \label{eq:K-kernel}
\end{align}
Although the displayed terms have poles, their specified differences
are regular.

\begin{lemma}[Holomorphy of the pole subtraction]
\label{lem:K-holomorphic}
For every fixed $|z|<1$, $\K(\alpha,\beta;z)$ extends holomorphically
to a neighborhood of $(0,0)$.
\end{lemma}

\begin{proof}
For $\Re\alpha>0$,
\begin{align*}
 \Li_{1+\alpha,1+\beta}(1,z)
 &=\sum_{n\ge1}\frac{z^n}{n^{1+\beta}}
   \zeta(1+\alpha,n+1),\\
 \Li_{1+\alpha,1+\beta}(1,z)
 -\frac{\Li_{1+\beta}(z)}{\alpha}
 &=\sum_{n\ge1}\frac{z^n}{n^{1+\beta}}
 \left(\zeta(1+\alpha,n+1)-\frac1\alpha\right).
\end{align*}
The parenthesis is holomorphic at $\alpha=0$.  On compact subsets of a
small bidisk it grows at most polynomially in $\log(n+1)$, whereas
$|z|^n$ decays exponentially.  The series is therefore locally
normally convergent and holomorphic in both variables.  The second
pole-subtracted term is symmetric, and the ordinary polylogarithm in
\eqref{eq:K-kernel} is already holomorphic in its order near two.
\end{proof}

Put
\begin{equation}
 G(\alpha)=\zeta(1+\alpha)-\frac1\alpha
 =\sum_{r\ge0}\frac{(-1)^r\gamma_r}{r!}\alpha^r.
 \label{eq:G-alpha}
\end{equation}

\begin{theorem}[All-index Stieltjes correlation generator]
\label{thm:stieltjes-correlation-generator}
For $|z|<1$ and $(\alpha,\beta)$ near $(0,0)$,
\begin{equation}
 \boxed{\displaystyle
 \sum_{n=0}^{\infty}
 \left(\zeta(1+\alpha,n+1)-\frac1\alpha\right)
 \left(\zeta(1+\beta,n+1)-\frac1\beta\right)z^n
 =\frac{G(\alpha)G(\beta)-\K(\alpha,\beta;z)}{1-z}.}
 \label{eq:parameter-correlation}
\end{equation}
Consequently, for all $r,s\ge0$,
\begin{equation}
 \boxed{\displaystyle
 \sum_{n=0}^{\infty}\gamma_r(n+1)\gamma_s(n+1)z^n
 =\frac{\gamma_r\gamma_s
 -(-1)^{r+s}r!s!\,[\alpha^r\beta^s]
       \K(\alpha,\beta;z)}{1-z}.}
 \label{eq:stieltjes-correlation}
\end{equation}
\end{theorem}

\begin{proof}
Apply \cref{thm:two-tail-generator} formally first in the half-plane
$\Re\alpha,\Re\beta>0$, with $p=1+\alpha$ and $q=1+\beta$.  Also use
the one-tail identity
\[
 \sum_{n\ge0}\zeta(1+\alpha,n+1)z^n
 =\frac{\zeta(1+\alpha)-\Li_{1+\alpha}(z)}{1-z}.
\]
Subtract $1/\alpha$ from the first factor and $1/\beta$ from the
second.  The zeta terms combine into $G(\alpha)G(\beta)$, while the
remaining polylogarithmic terms combine into $-\K$.  This proves
\eqref{eq:parameter-correlation} in a nonempty domain.  The holomorphic
extensions from \cref{lem:K-holomorphic} continue it to a neighborhood
of the origin.  Coefficient extraction using
\eqref{eq:stieltjes-laurent} gives
\eqref{eq:stieltjes-correlation}.
\end{proof}

\subsection{Digamma square and an anchored antiderivative}

Since $\gamma_0(n+1)=\gamma-H_n=-\psi(n+1)$, the lowest coefficient of
\cref{thm:stieltjes-correlation-generator} simplifies completely.

\begin{corollary}[Digamma-square generating function]
\label{cor:digamma-square}
For $|z|<1$,
\begin{equation}
 \boxed{\displaystyle
 \sum_{n=0}^{\infty}(\gamma-H_n)^2z^n
 =\sum_{n=0}^{\infty}\psi(n+1)^2z^n
 =\frac{(\gamma+\log(1-z))^2+\Li_2(z)}{1-z}.}
 \label{eq:digamma-square}
\end{equation}
\end{corollary}

\begin{proof}
The elementary generating functions
\[
 \sum_{n\ge0}H_nz^n=-\frac{\log(1-z)}{1-z},
 \qquad
 \sum_{n\ge0}H_n^2z^n
 =\frac{\Li_2(z)+\log^2(1-z)}{1-z}
\]
follow by summing the defining single and double harmonic sums.  Expand
$(\gamma-H_n)^2$ and combine them.
\end{proof}

Termwise integration provides the anchored primitive requested by the
factor $1/(n+1)$.

\begin{theorem}[Anchored digamma-square primitive]
\label{thm:digamma-square-primitive}
For $0\le z<1$,
\begin{align}
 \sum_{n=0}^{\infty}\frac{(\gamma-H_n)^2}{n+1}z^{n+1}
 ={}&\frac{\gamma^3-(\gamma+\log(1-z))^3}{3}
 -\Li_2(z)\log(1-z)-I(z),
 \label{eq:digamma-primitive}\\
 I(z):={}&\int_0^z\frac{\log^2(1-t)}t\dd t
 \notag\\
 ={}&2\zeta(3)+\log^2(1-z)\log z
 +2\log(1-z)\Li_2(1-z)-2\Li_3(1-z).
 \label{eq:I-closed}
\end{align}
In particular,
\begin{equation}
 \boxed{\displaystyle
 \sum_{n=0}^{\infty}
 \frac{(\gamma-H_n)^2}{(n+1)2^{n+1}}
 =\frac{\gamma^3-(\gamma-\log2)^3}{3}
 +\frac{\pi^2\log2}{12}-\frac{\log^32}{6}-\frac{\zeta(3)}4.}
 \label{eq:half-digamma-primitive}
\end{equation}
\end{theorem}

\begin{proof}
Integrate \eqref{eq:digamma-square} from $0$ to $z$.  The square term
is integrated by the substitution $u=\gamma+\log(1-t)$.  Integration
by parts gives
\[
 \int_0^z\frac{\Li_2(t)}{1-t}\dd t
 =-\Li_2(z)\log(1-z)-I(z).
\]
Differentiating the right side of \eqref{eq:I-closed} and checking its
zero value at $z=0$ proves the closed form.  At $z=1/2$, substitute the
classical evaluations
\[
 \Li_2(\tfrac12)=\frac{\pi^2}{12}-\frac{\log^22}{2},
 \qquad
 \Li_3(\tfrac12)=\frac78\zeta(3)-\frac{\pi^2\log2}{12}
                  +\frac{\log^32}{6}
\]
and simplify.
\end{proof}

More generally, integrating \eqref{eq:stieltjes-correlation} gives an
anchored antiderivative for every pair of indices:
\begin{equation}
 \sum_{n=0}^{\infty}
 \frac{\gamma_r(n+1)\gamma_s(n+1)}{n+1}z^{n+1}
 =\int_0^z
 \frac{\gamma_r\gamma_s-(-1)^{r+s}r!s!
 [\alpha^r\beta^s]\K(\alpha,\beta;t)}{1-t}\dd t.
 \label{eq:all-index-primitive}
\end{equation}
The point is not merely the integral notation: the integrand is a
finite coefficient of the explicitly specified depth-two kernel
\eqref{eq:K-kernel}, so order derivatives and anchored primitives are
controlled by one common analytic object.
 \section{An explicit \texorpdfstring{$(10,16)$}{(10,16)} reconstruction of \texorpdfstring{$\zeta(2)$}{zeta(2)}}
\label{sec:explicit-identity}

We record one fully expanded member of
\cref{thm:ordinary-reconstruction}.  The pair $(10,16)$ is the first
gap-six instance beyond the adjacent and gap-four order-ten cases.  Let
\begin{equation}
 L=14{,}612{,}936{,}162{,}752{,}272{,}720{,}199{,}680
 \label{eq:explicit-L}
\end{equation}
and define
\begin{equation}
 \mathscr P(x)=\sum_{r=0}^{15}c_rx^{2r},
 \label{eq:explicit-P}
\end{equation}
where the integer coefficients $c_r$ are listed in
\cref{tab:explicit-coefficients}.  Also put
\begin{align}
 \mathscr R(x)={}&
 1{,}552{,}076{,}939{,}623{,}430{,}612{,}844{,}544
 \notag\\
 &+2{,}244{,}965{,}091{,}528{,}662{,}071{,}640{,}064\,x^2
 \notag\\
 &+169{,}465{,}732{,}042{,}429{,}783{,}080{,}960\,x^4
 \notag\\
 &+2{,}907{,}927{,}622{,}440{,}764{,}571{,}648\,x^6.
 \label{eq:explicit-R}
\end{align}

\begin{theorem}[Explicit ordinary identity]
\label{thm:explicit-1016}
With $x_n=n+1/2$,
\begin{equation}
 \boxed{\displaystyle
 \sum_{n=0}^{\infty}
 \left\{
 135\,\mathscr P(x_n)
 [\zeta(10)-H_n^{(10)}][\zeta(16)-H_n^{(16)}]
 +\mathscr R(x_n)
 \right\}=L\zeta(2).}
 \label{eq:explicit-1016}
\end{equation}
The summand is $O(n^{-2})$.  Equivalently,
\begin{equation}
 \sum_{n=0}^{\infty}
 \left\{
 \frac{135\,\mathscr P(x_n)}{9!15!}
 \psi^{(9)}(n+1)\psi^{(15)}(n+1)
 +\mathscr R(x_n)
 \right\}=L\zeta(2).
 \label{eq:explicit-1016-polygamma}
\end{equation}
\end{theorem}

\begin{proof}
Take $b=8$ and $k=1$ in
\cref{thm:ordinary-reconstruction}.  Exact inversion of the rational
matrix $M_{5,8}$ gives
\[
 P_{8,1}(x)=\frac{\mathscr P(x)}
 {108{,}243{,}971{,}575{,}942{,}760{,}890{,}368}.
\]
The nonnegative asymptotic part \eqref{eq:Qbk} is
\begin{align*}
 Q_{8,1}(x)={}&-\frac{23127748662582496}{217749717276577245}
 -\frac{1238984692719613}{8064804343576935}x^2\\
 &-\frac{168349079333176}{14516647818438483}x^4
 -\frac{14443832747344}{72583239092192415}x^6.
\end{align*}
Multiplying
$\sum_n(P_{8,1}A_{10}A_{16}-Q_{8,1})=\zeta(2)$ by $L$, noting that
$L$ is $135$ times the denominator of $P_{8,1}$, gives
\eqref{eq:explicit-1016}.  Equation
\eqref{eq:explicit-1016-polygamma} follows from
\eqref{eq:polygamma-tail}.  The exact generator and verifier reconstruct
all displayed integers from \eqref{eq:source-M}, rather than checking a
pre-entered formula.
\end{proof}

The high-precision script evaluates the ordinary series directly and
agrees with $\zeta(2)$ after normalization to more than forty decimal
places.  This is recorded only as a diagnostic; the proof is the exact
matrix inversion and asymptotic subtraction above.
 \section{Audit findings, corrections, and integration consequences}
\label{sec:audit}

The source formulas used in the proof were rederived at the interfaces
where normalization mistakes are most likely: the omitted-binomial
convention in $M_{a,b}$, the factor two in the Bernoulli coordinates
$U_j=2B_{2b-2j}$, the retained-column reversal, and the scalar
multiplier in \eqref{eq:exact-multiplier}.  No mathematical
contradiction was found in those inherited formulas.  The following
status and presentation corrections should nevertheless be made when
this package is integrated.

\subsection{Research-status updates}

\begin{enumerate}[label=\textbf{A\arabic*.}]
\item The statement that the smaller-order-ten family remains open is
now obsolete.  It should be replaced by
\cref{thm:complete-order-ten}: every $M_{5,b}$ with $b\ge6$ is
nonsingular.

\item The exceptional congruence class
$b\equiv962\pmod{2048}$ in the preceding $2$-adic theorem was a
limitation of that congruence, not a candidate zero set.  The new finite
certificate includes every member of this class below $1920$, beginning
with $b=962$, and the inherited eventual theorem covers all members at
or above $1920$.

\item The broad unequal-order conjecture is still open.  Completion of
$p=10$ must not be described as a proof for all $p<q$.  The next fixed
smaller order is $p=12$, and the all-order problem requires control
uniform in $a$.
\end{enumerate}

\subsection{Certificate wording}

The residue vector is deliberately mixed: for $b<10$ it contains the
direct source determinant, while for $b\ge10$ it contains the normalized
five-square determinant.  Calling every entry ``the residue of
$\det M_{5,b}$'' would be inaccurate.  The exact definition is
\eqref{eq:proof-residue}; equivalence to source nonvanishing follows
from the nonzero multiplier.  The JSON metadata and integration notes
preserve this distinction.

A checksum alone is not a certificate.  The proof objects are the full
residue vector, the formulas generating each matrix, the denominator
argument, and the replay programs.  The SHA-256 value only detects
accidental or silent modification of that data.

\subsection{Analytic normalization}

The ordinary identities require the centered coordinate
$x_n=n+1/2$.  At this center the subtraction polynomial contains only
even powers and every such monomial has zero Hurwitz-zeta regularized
sum because $\zeta(-2j,1/2)=0$.  Replacing $x_n$ by $n+1$ without
recomputing the asymptotic polynomial would introduce nonzero Bernoulli
counterterms and invalidate \eqref{eq:ordinary-reconstruction}.
Likewise, omitting the constant term of $Q_{b,k}$ is not harmless:
$\zeta(0,1/2)=0$ is exactly what permits that constant to be subtracted
without changing the regularized value.

\subsection{Novelty and literature scope}

The inspected project already contains extensive work on harmonic-tail
reconstruction, Stieltjes correlations, and multiple polylogarithms.
This report claims the following relative contributions only:

\begin{itemize}
\item the exact completion of the $p=10$ fixed-order family;
\item the explicit all-coordinate ordinary identities that follow from
that completion;
\item the two-tail generator \eqref{eq:two-tail-generator} in the
normalization used here; and
\item the pole-subtracted kernel \eqref{eq:K-kernel} and its coefficient
formula for integer-shift Stieltjes products.
\end{itemize}

A targeted literature search found the standard multiple-polylogarithm,
Hurwitz-zeta, Bernoulli, and Stieltjes ingredients, but not these exact
combined formulas in this normalization.  This is not a comprehensive
priority search, and no global first-discovery claim is made.  The
article remains an unrefereed research draft.

\subsection{Recommended manuscript placement}

The determinant theorem belongs immediately after the current
fixed-order and eventual-rank material.  The ordinary reconstruction
and polylogarithmic generator can follow the harmonic-tail chapter as a
new all-coordinate consequence.  The Stieltjes kernel is logically
independent of the order-ten certificate and may instead be placed in
the shifted-Stieltjes chapter, with a cross-reference to the tail
generator from which it is obtained by singular-order regularization.
The supplied integration fragment and claim ledger use namespaced labels
to avoid collisions during that editorial split.
 \section{Further research questions and exact next targets}
\label{sec:research}

The results above suggest several concrete projects.  Each target is
phrased so that a successful answer has a checkable mathematical
endpoint rather than a larger table of experiments.

\begin{enumerate}[label=\textbf{R\arabic*.},leftmargin=*]
\item \textbf{Complete smaller order twelve.}
Apply the fixed-size reduction with $a=6$.  First compute a rigorous
eventual threshold and then certify the remaining finite interval with
one or several denominator-safe primes.  A stronger result would replace
the finite scan by a structural congruence or a sign argument for the
six-square determinant.

\item \textbf{Remove computation from the order-ten theorem.}
Find a recurrence, total-positivity statement, or arithmetic invariant
showing $\det\A^{[5]}\ne0$ for every $b\ge10$.  The modular vector in
this package is useful training data, but the desired endpoint is a
uniform proof explaining why cancellation cannot occur.

\item \textbf{Prove the unrestricted unequal-order conjecture.}
The eventual theorem controls $b$ for fixed $a$, but its threshold grows
with $a$.  Seek a bound on the lower Bernoulli-degree terms that is
uniform in the separation $b-a$, or a determinant factorization whose
nonvanishing is manifest for all $1\le a<b$.

\item \textbf{Closed inverse rows.}
The polynomials $P_{b,k}$ are explicit through matrix inversion but lack
a product formula.  Derive closed coefficient rows, perhaps through the
Pad\'e orthogonal polynomials and Christoffel--Darboux kernels.  A useful
answer should specify every coefficient and denominator directly, not
refer to an unevaluated inverse matrix.

\item \textbf{Minimal subtraction and accelerated identities.}
The present $Q_{b,k}$ leaves an $O(n^{-2})$ summand.  Subtract additional
negative asymptotic powers and compensate them by the known values
$\zeta(2j,1/2)$ to obtain arbitrarily rapid convergence.  Determine a
canonical acceleration with small rational coefficients and quantify
its arithmetic denominator growth.

\item \textbf{Boundary values of the two-tail generator.}
Continue \eqref{eq:two-tail-generator} through $z=1$ after the complete
Euler-operator subtraction.  Express the finite boundary values in a
minimal basis of multiple zeta values and order derivatives.  The
operator formula makes every candidate finite, but functional reduction
at high shifted order remains nontrivial.

\item \textbf{Rationally shifted and character-twisted tails.}
Replace $A_p(n)$ by $\zeta(p,n+a)$ with rational $a$, or by a Dirichlet
character tail.  Determine the analogue of the three-region numerator
in \eqref{eq:two-tail-generator}, the correct centered subtraction, and
the resulting coordinate space of Dirichlet $L$-values.

\item \textbf{Higher products.}
For three or more tails, partition the summation region by the minimum
index and derive depth-$d$ multiple-polylogarithm generators.  Then ask
for finite-rank reconstruction theorems and ordinary convergent
polygamma-product identities.  The main challenge is to organize the
ordered set partitions without hiding endpoint terms.

\item \textbf{Boundary Stieltjes correlations.}
Study the Abelian and finite-part limits of
\eqref{eq:stieltjes-correlation} as $z\to1$.  Identify exactly which
combinations of Stieltjes constants, multiple-zeta derivatives, and
powers of $\log(1-z)$ survive.  The pole-subtracted kernel should be
continued jointly in $(\alpha,\beta,z)$ rather than coefficient by
coefficient.

\item \textbf{Functional equations for the kernel.}
Investigate transformations $z\mapsto1-z$, $z\mapsto z/(z-1)$, and
roots of unity for $\K(\alpha,\beta;z)$.  A complete formula must retain
the order-derivative terms created by the pole subtractions.  Special
values could yield new exact sums of
$\gamma_r(n+1)\gamma_s(n+1)$ with algebraic weights.

\item \textbf{Higher anchored primitives and multiple Gamma functions.}
Integrate the all-index kernel repeatedly and determine when the result
reduces to classical polylogarithms, Barnes multiple Gamma functions,
or genuinely new depth.  The $r=s=0$ formula
\eqref{eq:digamma-primitive} is the first case of this hierarchy.

\item \textbf{Formal verification.}
Formalize the finite determinant reduction, the denominator-safe modular
certificate, and the min-index proof of
\eqref{eq:two-tail-generator}.  The finite-field portion is well suited
to reflection.  A later layer can formalize the analytic continuation
and centered regularization used in
\cref{thm:ordinary-reconstruction}.

\item \textbf{Independent literature and priority audit.}
Search multiple-zeta, Euler-sum, and Stieltjes-constant literature for
forms equivalent to \eqref{eq:K-kernel} under stuffle or iterated-integral
transformations.  Record equivalence maps explicitly.  This should be
completed before any publication-level novelty claim is made.
\end{enumerate}
 \appendix
\section{Certificate specification}
\label{app:certificate}

This appendix records the finite proof object independently of any
particular programming language.

\subsection{Input range and residue meaning}

The fixed parameters are
\[
 a=5,
 \qquad P=1{,}000{,}003,
 \qquad 6\le b\le1919.
\]
For $b=6,7,8,9$, the proof residue is the determinant of the direct
$b\times b$ matrix \eqref{eq:source-M}.  For $10\le b\le1919$, it is
the determinant of the $5\times5$ matrix \eqref{eq:retained-A} after
$U_j=2B_{2b-2j}$.  All operations are in $\F_P$.

\subsection{Canonical generation algorithm}

\begin{lstlisting}[language=Python,caption={Mathematical pseudocode for the residue vector.}]
P = 1000003
build factorials and inverse factorials modulo P through 3839
B[0] = 1
for m = 1,...,1919:
    B[2m] = ((2m-1)/2
              - sum_{j=1}^{m-1} binom(2m+1,2j) B[2j])/(2m+1)

for b = 6,...,1919:
    if b < 10:
        R[b] = det(M_{5,b}) mod P
    else:
        substitute U_j = 2 B[2b-2j], j=1,...,9
        R[b] = det(A^[5]) mod P
    assert R[b] != 0

serialize R[6],R[7],...,R[1919] as comma-separated decimal integers
compute SHA256 of those bytes
\end{lstlisting}

The expected digest is \eqref{eq:residue-digest}.  The JSON file also
contains the matrix ranges, maximum indices, count, full vector, and
sentinel entries.  The certificate contains no random choices.

\subsection{Two replay paths}

The primary replay uses the even recurrence
\eqref{eq:even-bernoulli-recurrence} and modular Gaussian elimination.
The independent replay uses the full Bernoulli recurrence
\[
 B_n=-\frac1{n+1}\sum_{k=0}^{n-1}\binom{n+1}{k}B_k
\]

directly for $0\le n\le3838$, verifies the vanishing of odd Bernoulli
residues beyond $B_1$, and computes determinants by subset expansion.
The two paths share only the mathematical formulas in the article and
the published JSON result.

\subsection{Sentinel values}

\begin{center}
\begin{tabular}{rll}
\toprule
$b$&matrix used&residue modulo $P$\\
\midrule
6&direct $M_{5,6}$&484126\\
8&direct $M_{5,8}$&572800\\
10&retained $\A^{[5]}$&668726\\
962&retained $\A^{[5]}$&989617\\
1919&retained $\A^{[5]}$&667530\\
\bottomrule
\end{tabular}
\end{center}

These values are convenient smoke tests, but validation requires the
entire vector and digest.
 \section{Coefficient table for the explicit identity}
\label{app:coefficients}

The coefficients in \eqref{eq:explicit-P} are as follows.  They are generated exactly by \nolinkurl{verification/generate_explicit_identity.py}.

\begin{longtable}{r@{\qquad}r}
\caption{Integer coefficients $c_r$ in $\mathscr P(x)=\sum_{r=0}^{15}c_rx^{2r}$.\label{tab:explicit-coefficients}}\\
\toprule
$r$&$c_r$\\
\midrule
\endfirsthead
\toprule
$r$&$c_r$\\
\midrule
\endhead
\bottomrule
\endfoot
0&22571931032959635\\
1&18694493678772529244\\
2&2408312382491948637488\\
3&111814818905053409281472\\
4&2361587601371307705179904\\
5&24757907369483659839458304\\
6&132829188131082664630743040\\
7&360605317517714764997902336\\
8&453247434535286514780798976\\
9&156976179448593567429427200\\
10&-151937057583031755486527488\\
11&-142696273664989191655653376\\
12&-40996749143564725151334400\\
13&-4714335942036423192870912\\
14&-209449736850990295941120\\
15&-2907927622440764571648\\
\end{longtable}

For a machine-readable record, \texttt{explicit\_10\_16\_zeta2.json} stores the denominator of $P_{8,1}$, all sixteen numerators, and the four rational coefficients of $Q_{8,1}$.
 \section{Reproducibility and evidential roles}
\label{app:reproducibility}

All paths below are relative to the package root.

\subsection{Exact proof checks}

\begin{lstlisting}[language=bash]
python3 verification/generate_order10_certificate.py
python3 verification/verify_order10_certificate.py
python3 verification/verify_order10_certificate_independent.py
python3 verification/verify_normalization_spots.py
python3 verification/generate_explicit_identity.py
python3 verification/verify_explicit_identity.py
\end{lstlisting}

The first two commands generate and replay the published modular vector.
The third is the independently written full replay.  The fourth checks
the retained/source multiplier at five boundary values.  The last two
regenerate the explicit $(10,16)$ identity from rational matrix
inversion and centered asymptotics.

\subsection{Numerical diagnostics}

\begin{lstlisting}[language=bash]
python3 verification/numerical_check_explicit_identity.py
python3 verification/numerical_check_generators.py
\end{lstlisting}

These scripts use high-precision numerical arithmetic to detect
transcription or branch errors.  They are not proof gates and are not
interval-arithmetic certificates.

\subsection{Complete replay and build}

\begin{lstlisting}[language=bash]
python3 verification/run_all.py
latexmk -pdf -interaction=nonstopmode -halt-on-error article.tex
\end{lstlisting}

The package has no compiled Python extension and no network dependency.
The exact scripts use only the Python standard library; the numerical
scripts require \texttt{mpmath}.  The PDF build requires a conventional
\TeX\ Live installation with \texttt{latexmk}.  The supplied
\texttt{Makefile} exposes the targets \texttt{verify}, \texttt{pdf},
and \texttt{all}.

\subsection{Claim classification}

\begin{center}
\begin{tabular}{p{0.22\linewidth}p{0.67\linewidth}}
\toprule
Class&Role\\
\midrule
Inherited theorem&Shifted-Hankel reduction, exact determinant multiplier,
and eventual threshold $B_5=1920$.\\
New exact theorem&Finite modular completion, complete order-ten rank,
ordinary reconstruction, two-tail generator, and Stieltjes kernel.\\
Exact computation&Finite-field determinant vector and rational
coefficient generation; both are replayable.\\
Numerical diagnostic&High-precision evaluation of displayed identities;
useful for error detection but not invoked in proofs.\\
Open problem&Every item in \cref{sec:research}; none is represented as
settled.\\
\bottomrule
\end{tabular}
\end{center}
 \begin{thebibliography}{99}

\bibitem{Repo}
V.~Reshetnikov and contributors,
\emph{ProveIt: Polylogarithms and their Arithmetic Bridges},
collective manuscript and incoming research, inspected at commit
\texttt{04bf92d507588a26abadbaee1beb1e7350c8bffb} (11 October 2026).
\url{https://github.com/VladimirReshetnikov/ProveIt/tree/04bf92d507588a26abadbaee1beb1e7350c8bffb/Analysis/Polylogarithms/docs/manuscript}.

\bibitem{BTD}
ProveIt research continuation,
\emph{Bernoulli Determinants, Cubic Tornheim Jets, and Canonical Digamma Integrals},
11 October 2026, at the repository snapshot in \cite{Repo}; especially
the shifted-Hankel and general small-determinant sections.

\bibitem{QSR}
ProveIt research continuation,
\emph{Harmonic-Tail Reconstruction, Quartic Trigamma Moments, and
Spectral Stieltjes Calculus}, 11 October 2026, at the repository snapshot
in \cite{Repo}; especially the complete order-eight theorem and the
effective eventual-rank theorem.

\bibitem{ATSD}
ProveIt research continuation,
\emph{Alternating Tails and Higher Spectral Defects}, 11 October 2026,
at the repository snapshot in \cite{Repo}; especially the five-square
order-ten contraction and the higher-precision Bernoulli congruence.

\bibitem{DLMF}
NIST Digital Library of Mathematical Functions,
Chapters 5, 24, and 25: Gamma, Bernoulli polynomials, Hurwitz zeta, and
polylogarithms.  Version 1.2.8, 15 September 2026.
\url{https://dlmf.nist.gov/}.

\bibitem{DilcherJiu2023}
K.~Dilcher and L.~Jiu,
\emph{Hankel determinants of shifted sequences of Bernoulli and Euler numbers},
Contributions to Discrete Mathematics \textbf{18} (2023), no.~2.
\url{https://doi.org/10.55016/ojs/cdm.v18i2.73995}.
Preprint: \url{https://arxiv.org/abs/2105.01880}.

\bibitem{DilcherJiu2022}
K.~Dilcher and L.~Jiu,
\emph{Hankel determinants of sequences related to Bernoulli and Euler polynomials},
International Journal of Number Theory \textbf{18} (2022).
\url{https://doi.org/10.1142/S179304212250021X}.
Preprint: \url{https://arxiv.org/abs/2007.09821}.

\bibitem{Goncharov}
A.~B.~Goncharov,
\emph{Multiple polylogarithms and mixed Tate motives},
preprint, 2001, arXiv:math/0103059.
\url{https://arxiv.org/abs/math/0103059}.

\bibitem{BorweinMPL}
J.~M.~Borwein, D.~M.~Bradley, D.~J.~Broadhurst, and P.~Lison\v{e}k,
\emph{Special values of multiple polylogarithms},
Transactions of the American Mathematical Society \textbf{353} (2001),
907--941.
\url{https://doi.org/10.1090/S0002-9947-00-02616-7}.

\bibitem{Coffey}
M.~W.~Coffey,
\emph{Series representations for the Stieltjes constants},
Rocky Mountain Journal of Mathematics \textbf{44} (2014), 443--477.
\url{https://doi.org/10.1216/RMJ-2014-44-2-443}.
Preprint: \url{https://arxiv.org/abs/0905.1111}.

\end{thebibliography}
 
\end{document}
